% ============================================================
% abstract.tex — Resumen y palabras clave
% Estado: Investigación en progreso (enfoque teórico-conceptual)
% ============================================================

\begin{abstract}
La fragmentación del conocimiento biomédico sobre plantas medicinales peruanas constituye una barrera significativa para la investigación farmacológica y etnobotánica regional. La evidencia científica sobre compuestos bioactivos, propiedades terapéuticas y aplicaciones farmacológicas de especies como \textit{Uncaria tomentosa}, \textit{Lepidium meyenii} y \textit{Croton lechleri} se encuentra dispersa en múltiples repositorios, bases de datos biomédicas y literatura heterogénea, dificultando su recuperación contextualizada y sistematización. Los sistemas de búsqueda convencionales, basados en concordancia léxica, presentan limitaciones críticas en dominios especializados, mientras que los modelos generativos de lenguaje exhiben fenómenos de alucinación que comprometen la fiabilidad de las respuestas en contextos biomédicos.

El presente trabajo presenta los fundamentos teóricos y el diseño conceptual de \sirca{}, una arquitectura semi-autónoma de descubrimiento científico biomédico que integra: (i)~un pipeline de Retrieval-Augmented Generation especializado para literatura biomédica; (ii)~memoria vectorial persistente basada en embeddings semánticos; (iii)~un agente semi-autónomo con estrategias híbridas de retrieval y reranking contextual; (iv)~un sistema de adquisición científica automatizada; y (v)~generación contextualizada con grounding documental. Este artículo se centra en la fundamentación teórica de la propuesta, el análisis del estado del arte en sistemas RAG biomédicos y la definición arquitectónica conceptual del sistema, estableciendo las bases para su implementación y evaluación experimental en fases posteriores de la investigación.
\end{abstract}

\begin{IEEEkeywords}
Retrieval-Augmented Generation, memoria vectorial, NLP biomédico, plantas medicinales peruanas, recuperación semántica, grounding documental, embeddings biomédicos, etnobotánica computacional.
\end{IEEEkeywords}
